CPE has been shown to work on small and medium scale language models.
It can only help as a safety technique if it also scales to current frontier models, and beyond.

Which way it should scale is not obvious a priori.
On one hand the number of parameters CPE acts upon is a smaller fraction of all parameters the bigger the model gets, which would make elicitation harder with scale.
On the other hand steering vectors, a method similar to CPE, work better on bigger models \citep{beaglehole2026}, which would suggest the opposite.
We therefore measured how the effect size of CPE scales with the size of the model.

A good way to evaluate this is the password-locked task.
New model organisms of password-locked sandbagging are relatively quick to train, and the percentage of recovered performance gives a natural scale for the effectiveness of CPE.

\subsection{The dense ladder}

We took models from the Qwen3 family at 0.6B, 1.7B, 4B, 8B, 14B and 32B parameters and trained each of them on MMLU-Pro \citep{wang2024} to answer correctly when a deployment tag is given.
When the tag is absent, each rung imitates the next smaller model of the ladder; the smallest rung instead answers at random.
Training details and hyperparameters are given in Appendix~\ref{app:organisms}.

We then trained CPE on the non-passworded questions, first on a fixed ratio of layers to target, and selected the factor with the largest uplift on the benchmark.
Afterwards, for all rungs whose reconstruction was not already perfect, we ran a small sweep over the layers to apply CPE at, to see whether other settings work better.
Table~\ref{tab:pwlock-dense} gives the results.

\begin{table}[htbp]
  \centering
  \small
  % Generated by exhibits/pwlock_tables/table_ladder.py -- do not edit by hand.
\begin{tabular}{lrrrrrl}
\toprule
 & \multicolumn{3}{c}{MMLU-Pro accuracy} & \multicolumn{2}{c}{Gap recovered (\%)} & \\
\cmidrule(lr){2-4}\cmidrule(lr){5-6}
Model & No pw. & Password & Gap & Fraction rule & Best band & Band \\
\midrule
0.6B & 0.190 & 0.413 & 0.223 & $81 \pm 15$\textsuperscript{\dag} & $90 \pm 9$ & [7,9]$\to$21 \\
1.7B & 0.221 & 0.466 & 0.245 & $89 \pm 17$\textsuperscript{\dag} & $84 \pm 8$ & [7,9]$\to$15 \\
4B & 0.316 & 0.539 & 0.223 & $87 \pm 9$ & $87 \pm 9$ & [8,12]$\to$20 \\
8B & 0.429 & 0.573 & 0.143 & $97 \pm 14$ & $97 \pm 14$ & [8,12]$\to$20 \\
14B & 0.357 & 0.657 & 0.301 & $94 \pm 6$ & $94 \pm 6$ & [9,13]$\to$22 \\
32B & 0.520 & 0.687 & 0.167 & $50 \pm 23$\textsuperscript{\dag} & $61 \pm 10$ & [23,30]$\to$43 \\
\bottomrule
\end{tabular}

  \caption{%
    Dense Qwen3 ladder. Recovery is $(\text{CPE} - \text{no pw.}) / (\text{pw.} - \text{no pw.})$.
    ``Fraction rule'' is the band from the depth-fraction rule $[0.225L, 0.325L] \to 0.55L$;
    ``best band'' is the winner of the band sweep, whose source/target layers are given in the last column.
    Cells marked \textsuperscript{\dag} were measured at $n=200$, the rest at $n=1200$.
    4B, 8B and 14B were never swept, so their two recovery columns are the same run.%
  }
  \label{tab:pwlock-dense}
\end{table}

CPE reconstructs almost all of the password-locked performance up to 14B, and then drops off at the 32B scale.
This suggests that CPE, at least with the methodology used here, works worse on bigger models.

\subsection{Mixture-of-experts rungs}

That conclusion rests on a single datapoint rather than a broad trend, so we sought further evidence from bigger models.
The Qwen3 dense models only reach the 32B scale, but there are larger Qwen models with a mixture-of-experts architecture.
We trained model organisms at the 35B, 122B and 397B scale.
Here again we swept the band CPE is applied at, though for the 397B model we swept only the two most promising settings, for time reasons.
Table~\ref{tab:pwlock-moe} gives the results.

\begin{table}[htbp]
  \centering
  \small
  % Generated by exhibits/pwlock_tables/table_ladder.py -- do not edit by hand.
\begin{tabular}{lrrrrrl}
\toprule
 & \multicolumn{3}{c}{MMLU-Pro accuracy} & \multicolumn{2}{c}{Gap recovered (\%)} & \\
\cmidrule(lr){2-4}\cmidrule(lr){5-6}
Model & No pw. & Password & Gap & Fraction rule & Best band & Band \\
\midrule
35B-A3B & 0.529 & 0.699 & 0.170 & $49 \pm 10$ & $63 \pm 10$ & [13,15]$\to$23 \\
122B-A10B & 0.570 & 0.785 & 0.215 & $7 \pm 9$ & $13 \pm 9$ & [10,16]$\to$36 \\
397B-A17B & 0.637 & 0.849 & 0.212 & $-4 \pm 22$\textsuperscript{\dag} & $0 \pm 9$ & [20,22]$\to$35 \\
\bottomrule
\end{tabular}

  \caption{%
    Qwen3.5 mixture-of-experts ladder, same columns and protocol as Table~\ref{tab:pwlock-dense}.%
  }
  \label{tab:pwlock-moe}
\end{table}

The trend continues, and CPE's ability to unlock the password-locked sandbagging has reached zero for the 397B model.

The differing architectures make it unclear whether the model sizes can usefully be compared with each other.
Figure~\ref{fig:pwlock-scaling} nevertheless gives a visual representation of CPE's apparent scaling behaviour.

\begin{figure}[htbp]
  \centering
  \includegraphics[width=0.92\textwidth]{pwlock_scaling/pwlock_scaling.png}
  \caption{%
    Recovered fraction of the password gap against model size, best band per rung.
    Dense rungs hold between 0.84 and 0.97 through 14B, and recovery falls away above that scale:
    0.61 at the dense 32B, and 0.63, 0.13, 0.00 across the MoE 35B, 122B and 397B.
    Error bars are $\pm 1$ standard error.%
  }
  \label{fig:pwlock-scaling}
\end{figure}

\subsection{Is it a steering-strength artifact?}

While this paints a pessimistic picture of CPE's scaling behaviour, it is possible that some aspect of how we apply CPE has to be scaled with model size.
A natural candidate is the steering strength $s$.
We did not have the time to conduct a sweep over different values of $s$ for the large models.
Instead we varied $s$ for the CPE LoRAs that were trained at $s=1$, and found that their performance degrades when $s$ is changed upwards or downwards (Figure~\ref{fig:pwlock-strength}).

\begin{figure}[htbp]
  \centering
  \includegraphics[width=0.58\textwidth]{pwlock_strength/pwlock_strength.png}
  \caption{%
    Recovery against steering strength $s$ for the two swept MoE rungs, re-scoring the selected factor at
    coefficients other than the default $s=1$ it was trained at ($n=600$).
    Both curves peak at $s=1$, so the weak recovery at scale is not a strength-calibration artifact.%
  }
  \label{fig:pwlock-strength}
\end{figure}

This is only one specific task, and it is not clear that CPE would show these scaling trends across all of its possible applications.
It does suggest that some part of how CPE is done has to be adjusted if we want to apply it to bigger models.
